% =====================================================================
%  Выводы по лабораторной работе.
% =====================================================================

\labpart{Выводы}

В ходе выполнения лабораторной работы был освоен механизм списков
контроля доступа: разработанная программа настраивает права доступа к
каталогам и файлам для всех предусмотренных заданием классов субъектов
(владелец, группа, остальные, все, администратор) и экспериментально
подтверждает их корректность прямыми проверками от имени каждого
пользователя. Практически показано, что классических прав \texttt{rwx}
достаточно не всегда -- выдать права ровно одной именованной группе можно
только через ACL, добавляемый утилитой \texttt{setfacl}.

Полученные результаты подтвердили и семантическую составляющую
теоретической части задания: доступ определяется единственной, наиболее
подходящей ACE, а для владельца объекта именованная запись группы не
учитывается вовсе, даже если он входит в эту группу. С прагматической
точки зрения работа показала практическую границу применимости
разрешения на выполнение -- оно защищает только скомпилированные
программы, а для исполняемых скриптов бесполезно без права на чтение,
поскольку их читает интерпретатор, а не ядро.
